Randomness is often treated as either ontic indeterminacy or epistemic ignorance. We develop a third, operational account within Six Birds: layer-relative randomness is the predictive residue of non-closure under a chosen packaging and budget. For a micro process $X_t$ and a packaging $\Pi : X \to Y$, we define the micro-closure deficit at scale $\tau$ by $\mathsf{CD}_\tau(\Pi)=I(X_t; Y_{t+\tau} \mid Y_t)$, the information about the next packaged state that remains hidden inside the current macro-object. This yields the exact decomposition $H(Y_{t+\tau}\mid Y_t)=H(Y_{t+\tau}\mid X_t)+\mathsf{CD}_\tau(\Pi)$, separating intrinsic substrate uncertainty from randomness introduced by discarded distinctions. We connect this exact quantity to computable diagnostics native to earlier Six Birds work, including route mismatch for coarse-grained Markov dynamics and predictive log-loss gaps under limited-memory models. In controlled Markov benchmarks, closure deficit vanishes on exactly lumpable partitions and rises with within-fiber heterogeneity, while route mismatch tracks it closely. In budgeted prediction, increasing memory order buys back hidden distinctions and lowers held-out log loss. In toy hashing, one-wayness and random-lookingness emerge as the same packaging-and-budget phenomenon: high-entropy inputs obey $q/2^n$ inversion scaling, while low-entropy inputs collapse the effect. The result is a measurable account of randomness that unifies coarse-graining, limited prediction, and feasible one-wayness without treating randomness as primitive.
